\section{Results and their precise interpretation}
Write $\N=\{0,1,2,\ldots\}$ and let $\Z$ denote unrestricted integers. Domains are part of each theorem and are never silently interchanged. Every independent coordinate is counted; an arithmetic abbreviation does not introduce another quantified variable. Coefficient height is measured after full collection, whereas residual rows and live binary operations are different quantities.

\subsection*{Natural quartics over $\N$}
\begin{center}\small
\begin{tabular}{@{}lrrrrr@{}}\toprule
Endpoint&Variables&Degree&Rows&Terms&Height\\\midrule
Strict cofinal anchor&229&4&179&9184&83954634\\
Sparse cofinal anchor&230&4&180&6858&83954472\\
High-height projection&228&4&178&9866&2445412949744\\
\bottomrule\end{tabular}\end{center}
\subsection*{Direct quartics over $\Z$}
\begin{center}\small
\begin{tabular}{@{}lrrrrr@{}}\toprule
Endpoint&Variables&Degree&Rows&Terms&Height\\\midrule
Signed dimension&282&4&192&5491&128446\\
Signed grouped&287&4&197&4313&128446\\
Signed sparse&309&4&219&3922&128446\\
\bottomrule\end{tabular}\end{center}
\subsection*{Height one quartics over $\Z$}
\begin{center}\small
\begin{tabular}{@{}lrrrrr@{}}\toprule
Endpoint&Variables&Degree&Terms&Height&Constant\\\midrule
Height one compact&793&4&51982&1&1\\
Height one balanced&806&4&23719&1&1\\
Height one sparse&994&4&11744&1&1\\
\bottomrule\end{tabular}\end{center}
\subsection*{Eleven integer arithmetic graphs}
\begin{center}\small
\begin{tabular}{@{}lrrr@{}}\toprule
Endpoint&Variables&Degree&Operations\\\midrule
Sharp free cofactor&11&328&130\\
Simple free cofactor&11&340&129\\
Uncentered free cofactor&11&352&125\\
Uncentered weak cofactor&11&398&121\\
\bottomrule\end{tabular}\end{center}


The natural cofinal-anchor family uses a single bounded-error row both to bound actual Pell indices and to match them to the requested indices. Aligning the amplification exponent with $x$ is justified by strict beta-quotient bounds, not by substituting $x$ into a theorem proved only for a larger exponent. Its smaller 228-variable companion has coefficient height $2445412949744$, so it is not the preferred balanced equation.

The signed family has its own complete sign proof. Orientation of the Pell pairs is established before deriving the signs of the exposed encoded representatives and quotient endpoint. This is not an unrestricted-domain reinterpretation of the natural certificate. Its witnesses are reconstructed with the cofinal representative and anchor choice in an explicit noncircular order.

The height-one family optimizes a different axis. It has more variables than the direct signed equations; the gain is that every collected nonzero coefficient is $1$ or $-1$. Actual incidence support, including harmless signed cancellations, replaces a stronger global collision prohibition. The independent check expands the whole final polynomial, so this improvement does not rely on an unverified graph heuristic.

For the eleven-variable family, linear cofactor saturation and the sharp polynomial center lower degree while preserving a complete authentic fixed coefficient tuple. The four recommended graphs are mutually nondominated in the reported degree/cost pair. Their actual highest-degree coefficients are nonzero at that tuple, not merely for generic coefficient symbols. Lower-cost higher-degree predecessors are retained as comparison points, not misclassified as dominated.

\subsection{The exact logical statement}
For each fixed equation $P$ on its stated domain $D$, the new ordinary arguments and the identified inherited calibration give
\[
 (\exists \boldsymbol{x}\in D^n\;P(\boldsymbol{x})=0)
 \quad\Longleftrightarrow\quad
 \text{there exists a finite syntactic ZFC refutation}.
\]
Consequently the claimed object-theoretic statement is
\[
 \mathrm{ZFC}\vdash
 \bigl[\forall\boldsymbol{x}\in D^n\;P(\boldsymbol{x})\ne0\bigr]
 \Longleftrightarrow\Con(\mathrm{ZFC}).
\]
Neither side is asserted to be true by this report. The finite checkers prove no consistency statement and do not decide the zero sets. No fresh transitive Lean or Rocq kernel build, encoded object-ZFC derivation, global minimum, or literature-wide novelty claim is made.

The direct logical interfaces are the retained named ordinary statements
\begin{quote}\small
\nolinkurl{KEndpoint.c222_iff_zf_refutation}\\
\nolinkurl{ConstructibleChoice.zf_refutation_iff_zfc_refutation}.
\end{quote}
The scalar route retains its source34/native/compiler/refutation theorem and the specifically identified signed-rank and interpolation arguments. The supplied numerical and executable closure is local, but this is not a new self-contained reconstruction of every historical native/compiler theorem. Inherited Lean files document the boundary; their inclusion alone is not a formal certification of the new equations.
